% ============================================================
%  styles/macros.tex
%  Custom commands, environments, and reusable building blocks
% ============================================================

% ── TEXT UTILITIES ───────────────────────────────────────────

% Highlight a term on first use
\newcommand{\term}[1]{\textbf{\textit{#1}}}

% Small caps abbreviation
\newcommand{\abbr}[1]{\textsc{\lowercase{#1}}}

% Emphasis that also adds to index
\newcommand{\idxterm}[1]{\textbf{#1}\index{#1}}

% Inline code snippet (allows line breaks)
\newcommand{\code}[1]{\texttt{\small\textcolor{CodeKeyword}{\hyphenchar\font=`\-\relax #1}}}

% File path (allows line breaks)
\newcommand{\filepath}[1]{\texttt{\small\textit{\hyphenchar\font=`\-\relax #1}}}

% Package name
\newcommand{\pkg}[1]{\textsf{#1}}

% Versioned software reference: \software{Docker}{20.10.7}
\newcommand{\software}[2]{\textit{#1}~v#2}

% ── MATH UTILITIES ───────────────────────────────────────────

% Norm  ||x||
\newcommand{\norm}[1]{\left\lVert#1\right\rVert}

% Absolute value |x|
\newcommand{\abs}[1]{\left\lvert#1\right\rvert}

% Sets
\newcommand{\R}{\mathbb{R}}
\newcommand{\N}{\mathbb{N}}
\newcommand{\Z}{\mathbb{Z}}

% ── THEOREM ENVIRONMENTS ─────────────────────────────────────
\theoremstyle{definition}
\newtheorem{definition}{Definition}[chapter]
\newtheorem{example}{Example}[chapter]
\newtheorem{remark}{Remark}[chapter]

\theoremstyle{plain}
\newtheorem{theorem}{Theorem}[chapter]
\newtheorem{lemma}[theorem]{Lemma}
\newtheorem{proposition}[theorem]{Proposition}

% ── CALLOUT BOXES (tcolorbox) ────────────────────────────────

% \infobox{Title}{Body text}
\newcommand{\infobox}[2]{%
  \begin{tcolorbox}[
    title       = #1,
    colback     = InfoBlue,
    colframe    = InfoBlueBorder,
    coltitle    = black,
    fonttitle   = \bfseries\small,
    breakable,
    left=6pt, right=6pt, top=4pt, bottom=4pt
  ]
  \small #2
  \end{tcolorbox}%
}

% \warnbox{Title}{Body text}
\newcommand{\warnbox}[2]{%
  \begin{tcolorbox}[
    title       = {\faExclamationTriangle\hspace{4pt}#1},
    colback     = WarnAmber,
    colframe    = WarnAmberBorder,
    coltitle    = black,
    fonttitle   = \bfseries\small,
    breakable,
    left=6pt, right=6pt, top=4pt, bottom=4pt
  ]
  \small #2
  \end{tcolorbox}%
}

% \tipbox{Title}{Body text}
\newcommand{\tipbox}[2]{%
  \begin{tcolorbox}[
    title       = #1,
    colback     = TipPurple,
    colframe    = TipPurpleBorder,
    coltitle    = white,
    fonttitle   = \bfseries\small,
    breakable,
    left=6pt, right=6pt, top=4pt, bottom=4pt
  ]
  \small #2
  \end{tcolorbox}%
}

% ── CHAPTER EPIGRAPH STYLE ───────────────────────────────────
\setlength{\epigraphwidth}{0.6\textwidth}
\renewcommand{\epigraphflush}{flushright}
\renewcommand{\epigraphrule}{0.4pt}

% ── FIGURE HELPERS ───────────────────────────────────────────

% Standard figure: \fig[options]{file}{caption}{label}
% options default to width=\textwidth
\NewDocumentCommand{\fig}{ O{width=\textwidth} m m m }{%
  \begin{figure}[htbp]
    \centering
    \includegraphics[#1]{#2}
    \caption{#3}
    \label{fig:#4}
  \end{figure}%
}

% Wide figure spanning full text width
\NewDocumentCommand{\widefig}{ m m m }{%
  \begin{figure}[htbp]
    \centering
    \includegraphics[width=\linewidth]{#1}
    \caption{#2}
    \label{fig:#3}
  \end{figure}%
}

% ── TABLE HELPERS ────────────────────────────────────────────

% Set global table spacing for a modern look
\renewcommand{\arraystretch}{1.3}

% Colored header row: \tableheader{Col1 & Col2 & Col3}
\newcommand{\tableheader}[1]{%
  \rowcolor{TableHeaderBg}%
  \bfseries\color{PrimaryDark}#1%
}

% Alternating row colors environment
\newenvironment{prettytable}[1]{%
  \rowcolors{2}{TableRowOdd}{TableRowEven}%
  \begin{tabular}{#1}%
  \toprule
}{%
  \bottomrule
  \end{tabular}%
}

% ── TABULARX COLUMN TYPES ────────────────────────────────────
\newcolumntype{L}{>{\raggedright\arraybackslash}X}
\newcolumntype{C}{>{\centering\arraybackslash}X}
\newcolumntype{R}{>{\raggedleft\arraybackslash}X}
\newcolumntype{P}[1]{>{\raggedright\arraybackslash}p{#1}}

% Column type that allows breaking for code-like text in tables
\newcolumntype{K}{>{\raggedright\arraybackslash\ttfamily\small}X}


% ── REFERENCE SHORTCUTS ──────────────────────────────────────
\newcommand{\figref}[1]{\hyperref[fig:#1]{Figure~\ref*{fig:#1}}}
\newcommand{\tabref}[1]{\hyperref[tab:#1]{Table~\ref*{tab:#1}}}
\newcommand{\chapref}[1]{\hyperref[chap:#1]{Chapter~\ref*{chap:#1}}}
\newcommand{\secref}[1]{\hyperref[sec:#1]{Section~\ref*{sec:#1}}}
\newcommand{\appref}[1]{\hyperref[app:#1]{Appendix~\ref*{app:#1}}}
\newcommand{\algref}[1]{\hyperref[alg:#1]{Algorithm~\ref*{alg:#1}}}
\newcommand{\diagref}[1]{\hyperref[diag:#1]{Diagram~\ref*{diag:#1}}}
\newcommand{\lstref}[1]{\hyperref[lst:#1]{Listing~\ref*{lst:#1}}}

% ── DRAFT UTILITIES ──────────────────────────────────────────
% Comment out for final version:
% \newcommand{\draftonly}[1]{#1}          % Show in draft
% \newcommand{\draftonly}[1]{}            % Hide in final
\newcommand{\draftonly}[1]{%
  \textcolor{red}{\textit{[DRAFT: #1]}}%
}

% ── PAGE BREAK CONTROL ───────────────────────────────────────
\newcommand{\chapterbreak}{\cleardoublepage}

% ── LOGO GRID (used in sec-4-software-stack) ─────────────────

\newcommand{\logocard}[2]{%
  \begin{minipage}[c]{2.8cm}
    \centering
    \IfFileExists{images/tools/#1.png}{%
      \includegraphics[height=1.0cm, keepaspectratio]{images/tools/#1}%
    }{%
      \fcolorbox{PrimaryDark}{LightGrayBg}{%
        \parbox[c][1.0cm][c]{2.2cm}{\centering\tiny\color{MidGray}#2}%
      }%
    }\\[4pt]
    {\tiny\color{DarkGray}#2}
  \end{minipage}%
}

\newenvironment{logogrid}{%
  \vspace{6pt}%
  \begin{center}%
  \setlength{\tabcolsep}{6pt}%
}{%
  \end{center}%
  \vspace{4pt}%
}